\begin{itemize}
    % Control Methods
    \item \textbf{Control Methods}
    \begin{itemize}
        \item \textbf{Adaptive Sliding Mode Fault-tolerant Control for Attitude Tracking of Spacecraft With Actuator Faults} \cite{Yang2023AdaptiveFaults} Proposes an adaptive sliding mode fault-tolerant control strategy using a two-phase power reaching law, ensuring robust attitude tracking in the presence of actuator faults.
        \item \textbf{Fault tolerant small satellite attitude control using adaptive non-singular terminal sliding mode} \cite{Cao2013FaultMode} Develops a sliding mode control approach with non-singular terminal properties to handle actuator faults in small satellite attitude control.
        \item \textbf{Indirect Robust Adaptive Fault-Tolerant Control for Attitude} \cite{Cai2008IndirectSpacecraft} Presents an indirect adaptive fault-tolerant control method designed for spacecraft attitude tracking in uncertain environments.
        \item \textbf{A Modified Robust Adaptive Fault Compensation Design for Spacecraft with Guaranteed Transient Performance} \cite{Yao2023APerformance} Introduces a robust adaptive control framework that ensures transient performance guarantees under actuator faults.
        \item \textbf{Adaptive Saturated Fault-Tolerant Control for Spacecraft Rendezvous with Redundancy Thrusters} \cite{Xia2021AdaptiveThrusters} Implements a saturated control mechanism integrated with an optimal fault-tolerant control allocation strategy for spacecraft rendezvous.
        \item \textbf{Actuator failure compensation and attitude control for rigid satellite by adaptive control using quaternion feedback} \cite{Ma2014ActuatorFeedback} Uses quaternion-based adaptive control to compensate for actuator failures in spacecraft attitude control.
        \item \textbf{An Application of Adaptive Fault-Tolerant Control to Nano-Spacecraft} \cite{Biggs2014AnNano-Spacecraft} Investigates adaptive fault-tolerant control strategies for nanosatellites, ensuring robust operation under actuator failures.
    \end{itemize}

    % Thruster Allocation Methods
    \item \textbf{Thruster Allocation Methods}
    \begin{itemize}
        \item \textbf{A New Control Allocation Methodology based on the Pseudo Inverse along the Null Space} \cite{BozorgniaD2013ASpace} Proposes an optimized control allocation technique leveraging the pseudo-inverse method within the null space.
        \item \textbf{Dual-Quaternion-Based Fault-Tolerant Control for Spacecraft Tracking With Finite-Time Convergence} \cite{Dong2017Dual-Quaternion-BasedConvergence} Uses a dual-quaternion framework for improved spacecraft tracking accuracy in the presence of faults.
        \item \textbf{Optimal Thrust Allocation Logic Design of Dynamic Positioning with Pseudo-Inverse Method} \cite{Yang2011OptimalMethod} Develops a pseudo-inverse method for optimal thruster allocation in dynamic positioning applications.
        \item \textbf{Resolving actuator redundancy—optimal control vs. control allocation} \cite{Harkegard2005ResolvingAllocation} Compares optimal control and control allocation methods for managing actuator redundancy.
        \item \textbf{Modified Pseudoinverse Redistribution Methods for Redundant Controls Allocation} \cite{Jin2005ModifiedAllocation} Proposes refinements to pseudo-inverse-based control allocation techniques for redundancy optimization.
        \item \textbf{Design of Thruster Control Allocation for a Cubesat} \cite{ElisaCAPELLOHyeongjunPARK2023POLITECNICOVITA} Examines control allocation methodologies tailored for CubeSat applications.
        \item \textbf{Lyapunov Stability Analysis of Daisy Chain Control Allocation} \cite{Buffington1996LyapunovAllocation} Analyzes daisy chain-based thruster allocation methods using Lyapunov stability theory.
        \item \textbf{Fault tolerant control design using adaptive control allocation based on the pseudo inverse along the null space} \cite{Tohidi2016FaultSpace} Introduces an adaptive control allocation framework to enhance fault tolerance in overactuated spacecraft.
        \item \textbf{Constrained Control Allocation} \cite{Durham1993ConstrainedAllocation} Investigates constraint-aware control allocation methods for aerospace applications.
        \item \textbf{CONSTRAINED CONTROL ALLOCATION FOR SYSTEMS WITH REDUNDANT CONTROL EFFECTORS} \cite{Bordignon1971ConstrainedI} Explores strategies for optimizing control allocation under redundancy constraints.
    \end{itemize}

    % Behavior Trees & System-Level Approaches
    \item \textbf{Behavior Trees \& System-Level Approaches}
    \begin{itemize}
        \item \textbf{Behavior Trees in Robot Control Systems} \cite{Ogren2022BehaviorSystems} Discusses behavior trees as a decision-making framework applicable to spacecraft autonomy and fault tolerance.
        \item \textbf{Bibliographical review on reconfigurable fault-tolerant control systems} \cite{Zhang2008BibliographicalSystems} Surveys fault-tolerant control techniques, comparing active and passive approaches in aerospace applications.
        \item \textbf{A Review on Recent Development of Spacecraft Attitude Fault Tolerant Control System} \cite{Yin2016ASystem} Provides an overview of recent advances in spacecraft fault-tolerant control systems.
        \item \textbf{Fault-tolerant control systems: A comparative study between active and passive approaches} \cite{Jiang2012Fault-tolerantApproaches} Evaluates active vs. passive fault-tolerant control approaches, detailing their benefits and limitations.
    \end{itemize}

    % Experimental & Simulation-Based Studies
    \item \textbf{Experimental \& Simulation-Based Studies}
    \begin{itemize}
        \item \textbf{Nanosatellite Air Bearing Tests of Fault-Tolerant Sliding-Mode Attitude Control with Unscented Kalman Filter} \cite{PostNanosatelliteFilter} Experiments with a fault-tolerant sliding mode control system for nanosatellites, using an air-bearing platform for validation.
        \item \textbf{Simulation and Implementation of Fault Tolerant Controllers for Attitude Control System of Spacecraft Test Bench} \cite{LeonSernaSimulationBench} Evaluates spacecraft fault-tolerant control strategies in a test bench environment.
        \item \textbf{Real-Time Implementation of Fault-Tolerant Control Systems with Performance Optimization} \cite{Yin2013Real-TimeOptimization} Investigates real-time performance considerations in fault-tolerant spacecraft control.
        \item \textbf{Development and testing of algorithms for optimal thruster command distribution during MTG orbital manoeuvres} \cite{DevelopmentManoeuvres} Presents optimization techniques for spacecraft thruster command distribution.
    \end{itemize}
    \item \textbf{TBD}
    \begin{itemize}
    % Fault Detection and Isolation (FDI) for Thrusters
    \item \textbf{Robust FDI applied to thruster faults of a satellite system} \cite{Patton2010RobustSystem} This paper explores fault detection and isolation (FDI) techniques for thruster failures in satellite systems. It evaluates Monte Carlo simulations and robust estimation methods to enhance reliability in fault detection, aiming to improve satellite operational safety.

    % Adaptive Fault-Tolerant Control
    \item \textbf{Adaptive Fault-Tolerant Spacecraft Attitude Control Design With Transient Response Control} \cite{Bustan2014AdaptiveControl} This study presents an adaptive fault-tolerant attitude control system with transient response optimization. The approach integrates backstepping and multiplicative fault modeling to handle actuator failures while maintaining stability and performance.

    % Control Allocation Methods
    \item \textbf{Control allocation—A survey} \cite{Johansen2013ControlSurvey} A comprehensive survey of control allocation techniques, including pseudo-inverse methods, optimization-based strategies, and constraint handling approaches. The paper discusses applications in aerospace, marine, and automotive control systems.

    % Passive Fault-Tolerant Control with Sliding Mode
    \item \textbf{Passive fault-tolerant sliding mode attitude control for flexible spacecraft with faulty thrusters} \cite{Mirshams2014PassiveThrusters} This work investigates passive fault-tolerant control for flexible spacecraft using adaptive sliding mode control. The approach compensates for actuator faults without requiring fault detection, enhancing robustness against unexpected failures.
\end{itemize}

\end{itemize}


\begin{itemize}
    \item \textbf{Random notes}
    \item Real time configurative control
    \item Variations of MPC superior to simpler path planning
    \item Typically, AFTCS can be divided into four sub-systems: (1)
         a reconfigurable controller, (2) a FDD scheme, (3) a controller
         reconfiguration mechanism, and (4) a command/reference
         governor.\cite{Zhang2008BibliographicalSystems}

 \item It should be noted that active FTC design incorporates FDD, 
        a controller reconfiguration mechanism, and a reconfigurable 
        feedback controller, while passive FTC design only consists of 
        a fixed fault tolerant controller design \cite{Yin2016ASystem}

        \item NASA’s Earth orbiting Lewis Spacecraft was launched on 
        August 23, 1997. It lost contact with the spacecraft on August 
        26, and then reentered the atmosphere and was destroyed on 
        September 28 \cite{Yin2016ASystem}

        \item Section III A in \cite{Yin2016ASystem} talked about past method for dealing with failures (non control, safe-mode -> manual plan)

        \item However, 
        the most applied allocation scheme is static allocation by using 
        the pseudo-inverse value of actuator distribution matrix [88]. 
        However, this approach is not an effective method to fully use 
        the remaining control power. As a result, several effective 
        dynamic control allocation schemes were proposed for 
        overactuated control system [89, 90]. \cite{Yin2016ASystem}

        \item \cite{Yang2011OptimalMethod} Uses PI with cost function combined using Lagrangian. (very cool)

        \item \cite{DevelopmentManoeuvres} "there is no difference in control authority between Simplex and Karmarkar, which have a slightly
        better control authority than the Pseudoinverse method"

        \item \cite{Johansen2013ControlSurvey} Mentions that PI is the "common" way to solve the allocation problem.
\end{itemize}